\caption{\textbf{Physical identification and learned residual forces improve equilibrium-shape prediction.}
(\textbf{A}) The conical, two-link, nine-DOF arm rendered in the learned-residual equilibrium, with a schematic mount and measured markers $M_1$--$M_4$ in gold. Gray spokes visualize the modeled backbone-to-marker offsets; wine crosses mark the predicted marker centres. The first three markers are offset by 25, 21, and 20~mm from the backbone, and the fourth extends 8~mm beyond the tip. The rendered surface is a model prediction; the experiment measures four marker positions.
(\textbf{B}) Initial-parameter, identified-parameter, and learned-residual predictions for single-tendon and combined-tendon inputs. Thick lines show backbones, fine spokes their modeled marker offsets, crosses predicted marker centres, and gold circles measurements. The faint surface shows the learned-residual prediction. Each configuration is selected by its proximity to the median identified-model RMS marker error within its input family; the tendon tensions are $(4.62,0)$, $(0,4.62)$, and $(4.62,6.15)$~N. All three views share camera and scale.
(\textbf{C}) RMS Euclidean marker-position error over the same 22 configurations and four markers for all stages: initial parameters, identified parameters, learned residuals, and independently fitted residuals as a reference. The residual model uses the rounded identified parameters. This pooled evaluation includes parameter-fitting and neural-network-training configurations and is not a held-out score.}
